% paper4_methods.tex -- Materials and Methods of Paper 4, input by
% papers/paper4_tillage.tex.
\section{Materials and Methods}

\subsection{Treatments and design}
Two tillage systems, conventional tillage (CT) and no-till (NT), were combined
factorially with two residue treatments, residues removed (R0) and residues
retained on the surface (R1), giving four combinations. The combinations were
arranged in randomised complete blocks\index{randomised complete block design} with \res{design-blocks} blocks at each of the three
sites and maintained for five years \parencite{gomez1984}. The data are held in
\texttt{data/trials/t4\_tillage.tsv}.

\subsection{Measurements}
Topsoil to \SI{0.20}{\metre} was sampled at the start and at the end of the
trial, and soil organic carbon was determined by wet oxidation\index{soil organic carbon!wet oxidation}
\parencite{walkley1934}. The change in soil organic carbon of a plot is
\begin{equation}
  \Delta\mathrm{SOC} = \mathrm{SOC}_{\mathrm{final}} - \mathrm{SOC}_{\mathrm{initial}}
  \quad (\si{\gram\per\kilo\gram}).
  \label{eq:dsoc}
\end{equation}
Maize grain yield in the final year is expressed in \si{\tonne\per\hectare} on a
dry-matter basis.

\subsection{Statistical model}
At each site, the response $y_{ijk}$ of tillage $i$ and residue treatment $j$ in
block $k$ is described by
\begin{equation}
  y_{ijk} = \mu + \beta_k + T_i + R_j + (TR)_{ij} + \varepsilon_{ijk},
  \label{eq:factorial}
\end{equation}
where $T_i$ and $R_j$ are the main effects of tillage and residues and
$(TR)_{ij}$ their interaction. A significant interaction indicates that the
effect of tillage depends on whether residues are retained.
